\documentclass[sigconf,natbib=true,anonymous=true]{acmart}

\AtBeginDocument{%
  \providecommand\BibTeX{{Bib\TeX}}}

\begin{document}

\title{Feature-Wise Contrastive Learning for Unsupervised Feature Selection}

\begin{abstract}
Unsupervised feature selection is important for high-dimensional data analysis,
yet many existing methods rely on graph construction, spectral objectives, or
reconstruction-based criteria that may be difficult to adapt across domains. We
propose a contrastive unsupervised feature selection method that treats each
feature as an instance and learns feature representations from multiple
perturbed views of its sample-profile vector. After training, features are
ranked using a simple embedding-based saliency score and evaluated by downstream
clustering performance. Experiments on benchmark datasets from multiple domains
show that the proposed method outperforms strong classical and neural baselines
on the selected evaluation suite. Additional analysis shows that the learned
representations contain meaningful ranking signal beyond a single fixed readout,
suggesting that feature-wise contrastive learning is a promising direction for
unsupervised feature selection.
\end{abstract}

\begin{CCSXML}
<ccs2012>
 <concept>
  <concept_id>10010147.10010257.10010293.10010294</concept_id>
  <concept_desc>Computing methodologies~Neural networks</concept_desc>
  <concept_significance>500</concept_significance>
 </concept>
 <concept>
  <concept_id>10002951.10003260.10003282</concept_id>
  <concept_desc>Information systems~Data mining</concept_desc>
  <concept_significance>300</concept_significance>
 </concept>
 <concept>
  <concept_id>10002951.10003260.10003283</concept_id>
  <concept_desc>Information systems~Clustering</concept_desc>
  <concept_significance>300</concept_significance>
 </concept>
</ccs2012>
\end{CCSXML}

\ccsdesc[500]{Computing methodologies~Neural networks}
\ccsdesc[300]{Information systems~Data mining}
\ccsdesc[300]{Information systems~Clustering}

\keywords{unsupervised feature selection, contrastive learning, feature ranking, clustering}

\maketitle

\section{Introduction}
High-dimensional datasets often contain many noisy, redundant, or weakly
informative features, which can degrade downstream clustering and exploratory
analysis. This challenge is especially common in biomedical, text, and
image-derived data, where the number of features can be large relative to the
number of samples and only a subset of features may reflect the underlying
structure of the data. Unsupervised feature selection (UFS) addresses this
problem by identifying a compact subset of informative features without access
to labels, thereby improving interpretability, reducing redundancy, and
supporting more effective unsupervised learning.

Compared with supervised feature selection, UFS is inherently more difficult
because there is no label signal to directly indicate which features are
discriminative or task-relevant. In supervised settings, feature importance can
be evaluated against class labels, whereas in unsupervised settings this signal
is unavailable, forcing the selection process to rely on indirect structural
surrogates such as similarity, locality, or cluster geometry. As a result, UFS
must infer feature utility from the intrinsic organization of the data alone,
which makes it harder to distinguish informative features from noisy or
redundant ones and to account for interactions among features, especially in
multi-cluster high-dimensional data.

As a result, prior work has relied on alternative principles for identifying
informative features. Early methods such as Laplacian Score (LS) rank each
feature according to how well it preserves local neighborhood relations in the
data \cite{he2005laplacian}. Multi-Cluster Feature Selection (MCFS) extends
this idea by selecting features that better retain the separation among
multiple clusters rather than evaluating each feature in isolation
\cite{cai2010unsupervised}. Other methods attempt to recover discriminative
structure without labels. For example, Nonnegative Discriminative Feature
Selection (NDFS) jointly estimates latent cluster assignments and the feature
subset that best supports them \cite{li2012unsupervised}, while Unsupervised
Discriminative Feature Selection (UDFS) seeks features that make the data more
separable under a learned discriminative transformation even without ground
truth labels \cite{yang2011udfs}. More recent neural approaches such as
Concrete Autoencoder (CAE) select a small subset of features that can be used
to reconstruct the full input \cite{balin2019concrete}, while LS-CAE augments
this reconstruction objective with a locality-preserving term so that the
selected features are encouraged to be both reconstructive and
structure-preserving \cite{shaham2022deep}. In contrast, we study UFS from a
feature-wise contrastive learning perspective, in which informative features are
identified through consistency across structured perturbations rather than
through neighborhood preservation, latent cluster recovery, or reconstruction
alone.

In this paper, we approach UFS through feature-wise contrastive learning.
Instead of treating samples as training instances, we invert the data matrix and
treat each feature as an instance represented by its sample-profile vector. For
each feature, we generate multiple structured perturbations and train the model
to produce consistent representations across views of the same feature while
distinguishing it from other features. After training, the learned embeddings
are converted into a ranking score, and the resulting ranking is used to select
feature subsets of different cardinalities. Under the standard clustering-based
UFS evaluation protocol \cite{shaham2022deep}, the proposed method achieves the best result on 8 out
of 10 benchmark datasets.

The main contributions of this paper are as follows:
\begin{itemize}
\item We propose a feature-wise contrastive learning method for unsupervised
feature selection that inverts the data matrix and learns representations of
features from structured perturbations of their sample profiles.
\item We introduce a simple embedding-based ranking strategy that produces a
single global feature ordering from one training run, enabling feature subset
selection at multiple cardinalities.
\item We show that the proposed method achieves strong clustering performance
against both classical and neural UFS baselines, obtaining the best result on 8
of 10 benchmark datasets under the standard evaluation protocol.
\end{itemize}

\section{Methodology}
Let $\mathbf{X} \in \mathbb{R}^{n \times d}$ denote a data matrix with $n$
samples and $d$ features. Instead of treating samples as training instances, we
invert the representation and treat each feature as an instance. Under this
feature-wise view, the $j$-th feature is represented by its sample-profile
vector $\mathbf{x}_j \in \mathbb{R}^{n}$, where $j \in \{1,\dots,d\}$.

For each feature $\mathbf{x}_j$, we construct four perturbed views
$\{\tilde{\mathbf{x}}_j^{(m)}\}_{m=1}^{4}$ using a four-tier masking regime. We
also construct a shuffled negative view of the same feature, denoted by
$\mathbf{x}_j^{-}$, by randomly permuting its sample order. An encoder $f_\theta$ maps
each feature view to a latent representation,
\[
\mathbf{h}_j = f_\theta(\mathbf{x}_j), \qquad
\tilde{\mathbf{h}}_j^{(m)} = f_\theta(\tilde{\mathbf{x}}_j^{(m)}),
\]
and a projector $g_\phi$ maps these latent representations to contrastive
embeddings,
\[
\mathbf{z}_j = g_\phi(\mathbf{h}_j), \qquad
\tilde{\mathbf{z}}_j^{(m)} = g_\phi(\tilde{\mathbf{h}}_j^{(m)}).
\]

After training, each feature is assigned a saliency score $s_j$. In our current
instantiation, we use the projector-side norm
\[
s_j = \|\mathbf{z}_j\|_2
\]
and rank features accordingly. The top-$k$ features form the selected subset
$\mathcal{S}_k$.




\subsection{Feature-Wise Formulation and Preprocessing}
We first standardize the input matrix $\mathbf{X} \in \mathbb{R}^{n \times d}$
across samples so that each feature has zero mean and unit variance. After
preprocessing, we invert the representation and work with the transposed matrix
$\mathbf{X}^{\top} \in \mathbb{R}^{d \times n}$, where each row corresponds to a
feature profile. In this feature-wise view, the $j$-th feature is treated as a
training instance and is represented by its sample-profile vector
$\mathbf{x}_j \in \mathbb{R}^{n}$. Following standard UFS practice
\cite{shaham2022deep}, the feature selection process is performed over the full
dataset, since the downstream evaluation is based on $k$-means clustering on
the full data. We then construct multiple structured perturbations of each
feature profile to define the contrastive learning task.\textbf{\textit{ bottom line: in ufs papers there is no train or test set, full datatset is scaled then the respective feature selection is performed and then k-means clsutering is performed on the full data and labels are used for measuring final k-meeans accuracy.  }}

\subsection{Perturbation Generation}
For each feature profile $\mathbf{x}_j$, we generate four positive views and one
structured negative view.

\begin{itemize}
\item Light mask view: a randomly selected $90\%$ subset of entries is retained.
\item Heavy mask view: a randomly selected $60\%$ subset of entries is retained.
\item Complementary view A: a randomly selected subset of entries is retained with keep ratio $0.50$.
\item Complementary view B: a second subset with the same keep ratio is constructed with controlled overlap ratio $0.10$ relative to view A.
\item Shuffled negative view: the sample order of $\mathbf{x}_j$ is randomly permuted.
\end{itemize}

The first four perturbed views are treated as positive views of the same
feature, while the shuffled view serves as an additional structured negative.
Beyond this explicit negative, other feature instances in the batch also act as
implicit negatives through the contrastive objective. Each view is then
processed by a shared encoder to obtain a feature-level latent representation.

\subsection{Encoder}
For each input view, the full set of feature instances is processed jointly by a
multi-head self-attention layer. Concretely, the view is represented as a
feature-token matrix in $\mathbb{R}^{d \times n}$, where $d$ is the number of
features, $n$ is the number of samples, each row corresponds to a feature
profile, and the sample-profile dimension serves as the token embedding
dimension. Self-attention allows each feature to attend to all other features in
the same view, producing a context-enriched representation before encoding.

The attention output is then passed to the encoder $f_\theta$, whose layer-wise
structure is summarized in Table~\ref{tab:arch}. The encoder maps the
feature-token matrix from $\mathbb{R}^{d \times n}$ to
$\mathbb{R}^{d \times h_e}$ and then to $\mathbb{R}^{d \times d_h}$ through a
two-stage multilayer perceptron, where $h_e$ denotes the encoder hidden
dimension and $d_h$ denotes the latent dimension. The resulting row vector
$\mathbf{h}_j \in \mathbb{R}^{d_h}$ is the encoder-side latent representation of
feature $j$. Analogously, each perturbed view produces a latent representation
$\tilde{\mathbf{h}}_j^{(m)} \in \mathbb{R}^{d_h}$, which is then passed to the
projector to obtain the contrastive embedding.

\subsection{Projector}
\label{sec:Projector}
We then apply a projector $g_\phi$ to map the encoder representation
$\mathbf{h}_j \in \mathbb{R}^{d_h}$ to a contrastive embedding
$\mathbf{z}_j \in \mathbb{R}^{d_z}$, where $h_p$ denotes the projector hidden
dimension and $d_z$ denotes the projector output dimension. The projector is
composed of two repeated projector blocks followed by a final linear
projection. Each projector block consists of a linear transformation, batch
normalization, and ReLU activation. As summarized in Table~\ref{tab:arch}, the
first projector block maps the encoder output from $\mathbb{R}^{d \times d_h}$
to $\mathbb{R}^{d \times h_p}$, and the second projector block further
transforms this intermediate representation from $\mathbb{R}^{d \times h_p}$ to
$\mathbb{R}^{d \times h_p}$.

A residual connection is inserted within the projector by adding the output of
the first projector block to the output of the second projector block. The
resulting representation is then passed through a final linear layer to obtain
the contrastive embedding in $\mathbb{R}^{d \times d_z}$. In our framework,
$\mathbf{z}_j$ is used for contrastive learning and downstream feature ranking,
while $\mathbf{h}_j$ serves as the pre-projection latent representation. These
projector embeddings define the space in which positive feature views are
aligned, negatives are separated, and the training objective is evaluated.



\begin{table}[t]
\centering
\caption{Architecture of the feature-wise encoder and projector. The encoder block consists of a linear transformation followed by batch normalization, LeakyReLU, and dropout, while each projector block consists of a linear transformation followed by batch normalization and ReLU activation.}
\label{tab:arch}
\small
\setlength{\tabcolsep}{5pt}
\renewcommand{\arraystretch}{1.1}
\resizebox{\columnwidth}{!}{
\begin{tabular}{|l|l|l|}
\hline
Layer & Input & Output \\
\hline
\multicolumn{3}{|c|}{\textbf{Input}} \\
\hline
Feature-token matrix & $\mathbf{X}^{(v)} \in \mathbb{R}^{d \times n}$ & $\mathbb{R}^{d \times n}$ \\
Multi-head self-attention over features & $\mathbb{R}^{d \times n}$ & $\mathbb{R}^{d \times n}$ \\
\hline
\multicolumn{3}{|c|}{\textbf{Encoder}} \\
\hline
Encoder block & $\mathbb{R}^{d \times n}$ & $\mathbb{R}^{d \times h_e}$ \\
Linear projection & $\mathbb{R}^{d \times h_e}$ & $\mathbb{R}^{d \times d_h}$ \\
\hline
\multicolumn{3}{|c|}{\textbf{Projector}} \\
\hline
Projector block 1 & $\mathbb{R}^{d \times d_h}$ & $\mathbb{R}^{d \times h_p}$ \\
Projector block 2 & $\mathbb{R}^{d \times h_p}$ & $\mathbb{R}^{d \times h_p}$ \\
Residual add & $\mathbb{R}^{d \times h_p}, \mathbb{R}^{d \times h_p}$ & $\mathbb{R}^{d \times h_p}$ \\
Linear projection & $\mathbb{R}^{d \times h_p}$ & $\mathbb{R}^{d \times d_z}$ \\
\hline
\end{tabular}
}
\end{table}


\subsection{Optimization Objective}
\paragraph{Contrastive Loss.}
Because the full set of features is processed jointly at each optimization step,
the contrastive objective is defined over all $d$ feature instances in the
dataset. Let $\mathbf{z}_j$ denote the anchor embedding of feature $j$, let
$\tilde{\mathbf{z}}_j^{(m)}$ denote the embedding of its $m$-th positive view,
and let $\mathbf{z}_j^{-}$ denote the embedding of its shuffled negative view.
Following the implementation, all embeddings are first $\ell_2$-normalized. For
each positive view $m \in \{1,\dots,4\}$, we optimize an InfoNCE-style
objective in which the matching pair
$(\mathbf{z}_j,\tilde{\mathbf{z}}_j^{(m)})$ is treated as the positive pair,
the embeddings of all other features act as implicit negatives, and the
shuffled view contributes an additional structured negative. The loss for view
$m$ is
\[
\mathcal{L}_{\mathrm{con}}^{(m)}
=
-\frac{1}{d}\sum_{j=1}^{d}
\log
\frac{
\exp\!\left(\langle \hat{\mathbf{z}}_j, \hat{\tilde{\mathbf{z}}}_j^{(m)} \rangle / \tau\right)
}{
\sum_{\ell=1}^{d}
\exp\!\left(\langle \hat{\mathbf{z}}_j, \hat{\tilde{\mathbf{z}}}_\ell^{(m)} \rangle / \tau\right)
+
\exp\!\left(\langle \hat{\mathbf{z}}_j, \hat{\mathbf{z}}_j^{-} \rangle / \tau\right)
}.
\]

\paragraph{Inter-Feature Decorrelation Regularizer.}
To discourage redundant feature representations, we regularize the normalized
anchor embeddings so that distinct features remain weakly correlated in the
projector space. Let $\hat{\mathbf{Z}} \in \mathbb{R}^{d \times d_z}$ denote the
matrix of $\ell_2$-normalized anchor embeddings for all $d$ features, where the
$j$-th row is $\hat{\mathbf{z}}_j \in \mathbb{R}^{d_z}$. We form the
inter-feature similarity matrix
\[
\mathbf{C} = \hat{\mathbf{Z}}\hat{\mathbf{Z}}^\top,
\]
whose entry
\[
C_{ij} = \left\langle \hat{\mathbf{z}}_i, \hat{\mathbf{z}}_j \right\rangle
\]
measures the cosine similarity between the normalized embeddings of features
$i$ and $j$.We then define the diversity loss as
\[
\mathcal{L}_{\mathrm{div}}
=
\frac{1}{d^2}
\left\|
\mathbf{C} - \mathbf{I}_d
\right\|_F^2
\]
where $\mathbf{I}_d$ is the $d \times d$ identity matrix. This regularizer suppresses off-diagonal similarities
between distinct feature embeddings while preserving unit self-similarity along
the diagonal, thereby encouraging a more diverse set of feature
representations.


\paragraph{Total Loss} The final loss function is summed up as:
    
\[
\mathcal{L}
=
\sum_{m=1}^{4}
\left(
\frac{1}{4}\mathcal{L}_{\mathrm{con}}^{(m)}
+
\lambda_{\mathrm{div}}\mathcal{L}_{\mathrm{div}}
\right).
\]
    


\subsection{Feature Selection Criterion}
After training, we use the projector-side representation as the downstream feature
representation and assign each feature a scalar saliency score
\[
s_j = \|\mathbf{z}_j\|_2.
\]
Features are ranked in descending order of $s_j$, and the top-$k$ features are
selected to form the subset $\mathcal{S}_k$.


\section{Experimental Setup}

\subsection{Evaluation Protocol}
For evaluation, we follow the standard unsupervised feature selection protocol used in LS-CAE~\cite{shaham2022deep} and prior UFS studies~\cite{li2012unsupervised}. Specifically, for each method we select feature subsets with cardinality
$k \in \{50,100,150,200,250,300\}$, run $k$-means on the selected features, and report the best clustering accuracy across the feature-count sweep. Following prior work, clustering quality is assessed using clustering accuracy.

All methods, including ours and the baselines, operate on the full dataset under
the same protocol. Labels are used only for evaluation.

\subsection{Datasets}
We evaluate on ten benchmark datasets spanning image, biomedical, text, and
mass-spectrometry domains, following the classical UFS benchmark suite
distributed through \texttt{scikit-feature}.\footnote{\url{https://jundongl.github.io/scikit-feature/datasets.html}}

\begin{table}[h]
  \caption{Summary of benchmark datasets.}
  \label{tab:datasets}
  \centering
  \begin{tabular}{lccc}
    \toprule
    Dataset & Samples & Features & Classes \\
    \midrule
    ALLAML     & 72   & 7129  & 2  \\
    ARCENE     & 200  & 10000 & 2  \\
    BASEHOCK   & 1993 & 4862  & 2  \\
    COIL20     & 1440 & 1024  & 20 \\
    LUNG       & 203  & 3312  & 5  \\
    NCI9       & 60   & 9712  & 9  \\
    PCMAC      & 1943 & 3289  & 2  \\
    PROSTATE   & 102  & 5966  & 2  \\
    RELATHE    & 1427 & 4322  & 2  \\
    WARPPIE10P & 210  & 2420  & 10 \\
    \bottomrule
  \end{tabular}
\end{table}



\subsection{Baselines}
We compare against representative classical and neural unsupervised feature
selection baselines. The classical baselines include Laplacian Score (LS)~\cite{he2005laplacian},
Multi-Cluster Feature Selection (MCFS)~\cite{cai2010unsupervised},
Nonnegative Discriminative Feature Selection (NDFS)~\cite{li2012unsupervised},
and Unsupervised Discriminative Feature Selection (UDFS)~\cite{yang2011udfs}.
These methods cover graph-based, spectral, and sparse discriminative UFS
formulations. We also compare against neural baselines, including Concrete
Autoencoder (CAE)~\cite{balin2019concrete} and LS-CAE~\cite{shaham2022deep}.

\section{Results and Analysis}
\subsection{Main Results}
Table~\ref{tab:main-results} reports the best clustering accuracy under the
standard feature-count sweep. Overall, ICL achieves the best result on 8 of the
10 datasets and the second-best result on RELATHE, where it is only 0.16 points
below the strongest baseline. The only dataset on which ICL does not rank among
the top two methods is COIL20, where LS-CAE attains the best performance.

ICL consistently outperforms both classical and neural baselines on most
datasets. Its gains over the strongest competing baseline are especially large
on PROSTATE (+11.38) and WARPPIE10P (+11.12), while it also improves on
ALLAML, ARCENE, BASEHOCK, LUNG, NCI9, and PCMAC. Compared with LS-CAE, ICL
performs better on 9 of the 10 datasets, indicating that the proposed inverted
contrastive formulation provides a strong and broadly effective alternative to
existing unsupervised feature selection methods under the standard
clustering-based protocol.


\begin{table*}[h]
  \caption{Best clustering accuracy (\%) under the standard feature-count sweep. Best is in bold and second best is underlined.}
  \label{tab:main-results}
  \centering
  \begin{tabular}{lccccccc}
    \toprule
    Dataset & LS & MCFS & NDFS & UDFS & CAE & LS-CAE & ICL \\
    \midrule
    ALLAML     & 63.61 & 70.76 & 67.08 & -- & 55.49 & \underline{70.83} & \textbf{73.33} \\
    ARCENE     & 63.88 & 63.90 & \underline{64.70} & -- & 62.43 & \underline{64.70} & \textbf{66.33} \\
    BASEHOCK   & 50.87 & 51.21 & 50.68 & -- & \underline{51.44} & 50.95 & \textbf{52.42} \\
    COIL20     & 56.31 & 60.78 & 57.86 & -- & \underline{64.07} & \textbf{64.76} & 61.89 \\
    LUNG       & 57.07 & \underline{59.26} & 54.53 & -- & 53.42 & 57.27 & \textbf{63.82} \\
    NCI9       & \underline{41.08} & 37.83 & -- & -- & 39.08 & 37.17 & \textbf{43.50} \\
    PCMAC      & 50.61 & 50.67 & 50.81 & -- & \underline{51.31} & 50.78 & \textbf{51.65} \\
    PROSTATE   & 57.21 & 57.84 & 56.08 & -- & \underline{60.78} & 59.85 & \textbf{72.16} \\
    RELATHE    & 54.23 & 54.94 & 54.57 & -- & \textbf{55.87} & 55.46 & \underline{55.71} \\
    WARPPIE10P & 37.43 & 31.57 & \underline{39.64} & -- & 28.36 & 32.57 & \textbf{50.76} \\
    \bottomrule
  \end{tabular}
\end{table*}


\textbf{\textit{i have added 3 candidate tables below that provides icl vs strongest baseline in 3 different temp, cross correlation combination. let me know which one to pick for the full table. all 3 combinations give 8/10 wins, just at different margins. the current main table uses \ref{tab:icl-vs-baseline-003-035}}}

\begin{table}[t]
\centering
\caption{Baseline vs. ICL for $(\tau,\lambda_{\mathrm{div}})=(0.03,0.30)$. This is neither the most balanced version neither the most "statement wins" version. its just the only other combination with a 8/10 table other than the other two}
\label{tab:icl-vs-baseline-003-030}
\resizebox{\columnwidth}{!}{%
\begin{tabular}{lccccc}
\toprule
Dataset & Best Baseline & Baseline Acc. & ICL $k$ & ICL Acc. & Gap \\
\midrule
ALLAML     & LS-CAE @ 300 & 70.83 & 100 & \textbf{73.12} & +2.29 \\
ARCENE     & NDFS @ 100   & 64.70 & 100 & \textbf{66.50} & +1.80 \\
BASEHOCK   & CAE @ 250    & 51.44 & 300 & \textbf{52.34} & +0.90 \\
COIL20     & LS-CAE @ 250 & \textbf{64.76} & 250 & 63.08 & -1.68 \\
LUNG       & MCFS @ 100   & 59.26 & 150 & \textbf{64.01} & +4.75 \\
NCI9       & LS @ 300     & 41.08 & 150 & \textbf{43.67} & +2.58 \\
PCMAC      & CAE @ 100    & 51.31 & 50  & \textbf{51.79} & +0.48 \\
PROSTATE   & CAE @ 50     & 60.78 & 50  & \textbf{68.09} & +7.30 \\
RELATHE    & CAE @ 250    & \textbf{55.87} & 50  & 55.68 & -0.20 \\
WARPPIE10P & NDFS @ 50    & 39.64 & 200 & \textbf{41.43} & +1.79 \\
\bottomrule
\end{tabular}%
}
\end{table}

\begin{table}[t]
\centering
\caption{Baseline vs. ICL for $(\tau,\lambda_{\mathrm{div}})=(0.03,0.35)$. statement wins in prostate and war(last row)  11 \%+ with similar margin of wins as the previous table in other datastets, only downside a bigger loss in coil20 probably the best combination}
\label{tab:icl-vs-baseline-003-035}
\resizebox{\columnwidth}{!}{%
\begin{tabular}{lccccc}
\toprule
Dataset & Best Baseline & Baseline Acc. & ICL $k$ & ICL Acc. & Gap \\
\midrule
ALLAML     & LS-CAE @ 300 & 70.83 & 300 & \textbf{73.33} & +2.50 \\
ARCENE     & NDFS @ 100   & 64.70 & 100 & \textbf{66.33} & +1.63 \\
BASEHOCK   & CAE @ 250    & 51.44 & 100 & \textbf{52.42} & +0.98 \\
COIL20     & LS-CAE @ 250 & \textbf{64.76} & 200 & 61.89 & -2.88 \\
LUNG       & MCFS @ 100   & 59.26 & 300 & \textbf{63.82} & +4.56 \\
NCI9       & LS @ 300     & 41.08 & 50  & \textbf{43.50} & +2.42 \\
PCMAC      & CAE @ 100    & 51.31 & 50  & \textbf{51.65} & +0.34 \\
PROSTATE   & CAE @ 50     & 60.78 & 50  & \textbf{72.16} & +11.37 \\
RELATHE    & CAE @ 250    & \textbf{55.87} & 50  & 55.71 & -0.16 \\
WARPPIE10P & NDFS @ 50    & 39.64 & 50  & \textbf{50.76} & +11.12 \\
\bottomrule
\end{tabular}%
}
\end{table}



\begin{table}[t]
\centering
\caption{Baseline vs. ICL for $(\tau,\lambda_{\mathrm{div}})=(0.05,0.25)$. The most balanced one, higher margins ins basehock and only combination winning in  coil20, lowest margine of loss in the two losing dataste, also prsotate and war(last row) gains are decent and balanced}
\label{tab:icl-vs-baseline-005-025}
\resizebox{\columnwidth}{!}{%
\begin{tabular}{lccccc}
\toprule
Dataset & Best Baseline & Baseline Acc. & ICL $k$ & ICL Acc. & Gap \\
\midrule
ALLAML     & LS-CAE @ 300 & 70.83 & 300 & \textbf{72.64} & +1.81 \\
ARCENE     & NDFS @ 100   & 64.70 & 50  & \textbf{66.00} & +1.30 \\
BASEHOCK   & CAE @ 250    & 51.44 & 150 & \textbf{53.40} & +1.96 \\
COIL20     & LS-CAE @ 250 & 64.76 & 300 & \textbf{67.26} & +2.50 \\
LUNG       & MCFS @ 100   & 59.26 & 300 & \textbf{61.72} & +2.46 \\
NCI9       & LS @ 300     & 41.08 & 100 & \textbf{46.92} & +5.83 \\
PCMAC      & CAE @ 100    & \textbf{51.31} & 150 & 50.93 & -0.38 \\
PROSTATE   & CAE @ 50     & 60.78 & 50  & \textbf{66.42} & +5.64 \\
RELATHE    & CAE @ 250    & \textbf{55.87} & 250 & 55.07 & -0.80 \\
WARPPIE10P & NDFS @ 50    & 39.64 & 50  & \textbf{46.79} & +7.14 \\
\bottomrule
\end{tabular}%
}
\end{table}
\textbf{these are not the final results, i still have sweeps left in projector embedding size and encoder embedding size, maybe masking percentage too. Final table by tomorrow}

\subsection{Ablation Study}
We evaluate the contribution of the main design components through a progressive
ablation. Starting from the full model, we remove the diversity regularizer, then
remove the attention layer, and finally replace the structured four-view masking
scheme with simple random masking. This produces a compact hierarchy from the
full method to a minimal contrastive baseline.
\begin{table*}[h]
  \caption{Ablation study using best clustering accuracy (\%). ``w/o Div'' removes the diversity loss, ``w/o Attn'' removes the attention layer, and ``1-View RandMask'' replaces the structured four-view masking with a single random-masked positive view.}
  \label{tab:ablation}
  \centering
  \begin{tabular}{lcccc}
    \toprule
    Dataset & Full & w/o Div & w/o Div + w/o Attn & w/o Div + w/o Attn + 1-View RandMask \\
    \midrule
    COIL20   & \textbf{65.94} & \underline{61.20} & 59.69 & 58.72 \\
    ALLAML   & \textbf{72.22} & \textbf{72.22} & \underline{72.01} & 68.82 \\
    PROSTATE & \underline{71.42} & 71.57 & 57.21 & \textbf{74.12} \\
    BASEHOCK & \underline{52.38} & \textbf{53.06} & 52.33 & 50.58 \\
    \bottomrule
  \end{tabular}
\end{table*}

Table~\ref{tab:ablation} shows that the full model is generally the most stable
configuration, with the strongest overall performance across datasets. Removing
attention consistently hurts COIL20 and causes a substantial drop on PROSTATE,
indicating that feature-wise interaction modeling is an important component of
the method. Replacing the structured four-view construction with a single random
masked view further degrades performance on COIL20, ALLAML, and BASEHOCK,
suggesting that the proposed perturbation design contributes meaningfully beyond
simple masking. The effects of the diversity term are smaller and dataset
dependent, providing modest gains in some cases while acting primarily as a
stabilizing regularizer rather than the sole source of improvement.

\subsection{Alternate Selection Criterion}
Our default feature ranking criterion uses the encoder-side embedding norm
$s_j=\|\mathbf{z}_j\|_2$. To assess whether the learned representations contain
additional ranking signals beyond this default choice, we compare it against a
small set of alternatives derived from the encoder and projector outputs. In
addition to \emph{$z$-norm}, we consider \emph{$h$-norm}, which ranks features by
the projector-space magnitude $\|\mathbf{h}_j\|_2$, and \emph{harmonic-$hz$},
which computes the harmonic mean of the min-max scaled $h$-norm and $z$-norm
scores. We also evaluate a contrastive consistency criterion defined as
\[
s_j^{\mathrm{cc}}
=
\frac{\mu_j^{(+)}-\mu_j^{(-)}}{\sigma_j^{(+)}+\epsilon},
\]
where $\mu_j^{(+)}$ is the mean anchor-positive cosine similarity across the
four positive views, $\mu_j^{(-)}$ is the anchor-negative cosine similarity, and
$\sigma_j^{(+)}$ is the standard deviation across positive-view similarities.
This score favors features that are consistently aligned with their positive
views while remaining distinct from the shuffled negative.

Table~\ref{tab:alt-score} shows that the effectiveness of these criteria is
dataset dependent. The default $h$-norm remains strongest on COIL20 and
PROSTATE, while the contrastive consistency criterion performs best on ALLAML.
Projector-space norm is particularly competitive and provides the best result on
BASEHOCK, making it the strongest alternative to the default ranking rule in our
study. However, no single criterion consistently dominates across all datasets.
This suggests that the learned representations do contain multiple useful
ranking signals, but changing the readout alone yields only limited gains; more
substantial improvements are likely to require a better alignment between the
contrastive task and the desired notion of feature utility.
\begin{table*}[h]
  \caption{Best clustering accuracy (\%) under alternative feature scoring criteria. We compare the default encoder-side norm against projector-space norm, a harmonic fusion of scaled norms, and the proposed instability-margin criterion. Best is in bold and second best is underlined.}
  \label{tab:alt-score}
  \centering
  \begin{tabular}{lcccc}
    \toprule
    Dataset & $h$-norm & $z$-norm & Harmonic-$hz$ & Instability Margin \\
    \midrule
    COIL20   & \textbf{65.94} & 62.42 & \underline{64.49} & 58.39 \\
    ALLAML   & 72.22 & \underline{72.78} & 72.22 & \textbf{73.06} \\
    PROSTATE & \textbf{71.42} & 59.26 & \underline{65.20} & 61.81 \\
    BASEHOCK & 52.38 & \textbf{55.58} & \underline{53.51} & 51.95 \\
    \bottomrule
  \end{tabular}
\end{table*}

\subsection{Hyperparameter Analysis}
\begin{table}[h]
  \caption{Best diversity-loss weight per dataset.}
  \label{tab:best-div}
  \centering
  \begin{tabular}{lccc}
    \toprule
    Dataset & Best $\lambda_{\mathrm{div}}$ & Best $k$ & ACC (\%) \\
    \midrule
    COIL20   & 0.07  & 300 & 65.94 \\
    ALLAML   & 0.001 & 100 & 72.22 \\
    PROSTATE & 0.03  & 50  & 72.55 \\
    BASEHOCK & 0.00  & 100 & 53.06 \\
    \bottomrule
  \end{tabular}
\end{table}
\begin{figure}[h]
  \centering
  \includegraphics[width=\linewidth]{plots/diversity_sweep.png}
  \caption{Sensitivity to diversity-loss weight. Starred markers indicate the best setting for each dataset.}
  \label{fig:diversity-sweep}
\end{figure}

\begin{table}[h]
  \caption{Best temperature per dataset.}
  \label{tab:best-temp}
  \centering
  \begin{tabular}{lccc}
    \toprule
    Dataset & Best $\tau$ & Best $k$ & ACC (\%) \\
    \midrule
    COIL20   & 0.06 & 200 & 66.47 \\
    ALLAML   & 0.19 & 50  & 72.71 \\
    PROSTATE & 0.08 & 50  & 71.52 \\
    BASEHOCK & 0.09 & 50  & 53.52 \\
    \bottomrule
  \end{tabular}
\end{table}
\begin{figure}[h]
  \centering
  \includegraphics[width=\linewidth]{plots/temperature_sweep.png}
  \caption{Sensitivity to contrastive temperature. Starred markers indicate the best setting for each dataset.}
  \label{fig:diversity-sweep}
\end{figure}


\section{Conclusion}
We introduced a feature-wise contrastive learning approach for unsupervised feature
selection. The method is simple, empirically competitive, and exhibits an
interesting two-stage behavior that separates early grouping from late feature
saliency formation. Future work will focus on designing task-aligned contrastive
objectives that better connect contrastive solvability to non-redundant and
clustering-useful feature subsets.

% Optional appendix: keep within the 4-page content limit if used.
% \appendix
% \section{Additional Experimental Details}

\bibliographystyle{ACM-Reference-Format}
\bibliography{references}

\end{document}
